\documentclass[10pt]{article}
\usepackage[T1]{fontenc}
\usepackage{lmodern,microtype,amsmath,amssymb,mathtools,array,tabularx,booktabs}
\usepackage[margin=0.85in]{geometry}
\usepackage{xurl,hyperref}
\hypersetup{colorlinks=true,urlcolor=blue,linkcolor=blue,pdftitle={Appendix A: corrections and independent reviews}}
\setlength{\parindent}{0pt}
\setlength{\parskip}{5pt}
\newcommand{\D}{\mathcal D}
\newcommand{\Two}{\mathbf 2}
\newcommand{\Tm}{\mathsf{Tm}}
\newcommand{\Fm}{\mathsf{Fm}}
\newcolumntype{Y}{>{\raggedright\arraybackslash}X}
\newcolumntype{L}[1]{>{\raggedright\arraybackslash}p{#1}}
\begin{document}
{\LARGE Appendix A: corrections and independent reviews}\par
{\large Syntax and deductively closed theories}\par
31 October 2026 \hfill Revision accompanying manuscript version 14

The two indicated corrections have been implemented. Separate agents applied
the hostile-review and mathematical-writing workflows; a third prepared the
actual six-section outline, and a fourth reconciled their reports and edited
both copies of the appendix. The new reviews identified no further necessary
repair. They support retaining the six sections and their dependencies.

This report concerns the standalone Appendix A, whose mathematical body is
Appendix B of \emph{Local connectedness in recursion categories and essential
incompleteness}. It records an author-side audit, not independent human peer
review or formal verification. The complete individual reports accompany this
summary in the repository's version 14 audit directory.

\section*{Actual six-section outline}
\renewcommand{\arraystretch}{1.16}
\begin{tabularx}{\textwidth}{@{}L{0.07\textwidth}L{0.31\textwidth}Y@{}}
\toprule
Section & Title & Function and dependencies\\\midrule
A.1 & Categorical hypotheses and notation
& Fixes the ambient category, coproduct assumptions, Boolean morphisms,
decidable variable equality, signatures, and parameter stability.\\
A.2 & Terms, formulas, and free variables
& Defines the initial syntax and list algebras; constructs list operations,
free-variable lists, and the sentence subobject. Uses A.1.\\
A.3 & Admissible substitution
& Constructs single-variable substitution and its admissibility subobject;
proves the identity and absent-variable laws. Uses A.1--A.2.\\
A.4 & Inference rules and deductively closed subobjects
& Defines the inference endofunctor in a slice and proves that the image of
its initial algebra is the least closed subobject. Uses regularity and finite
coproducts, independently of the first-order syntax.\\
A.5 & First-order theories
& Specifies logical schemes and rule morphisms; applies A.4 to A.2--A.3;
constructs sentence theories and their partial-identity representations.\\
A.6 & Propositional theories
& Applies A.4 to propositional syntax under weaker hypotheses, reusing the
three propositional schemes displayed in A.5.\\\bottomrule
\end{tabularx}

The principal result is Theorem A.5, the least closed image theorem.
Proposition A.3 supplies substitution for the first-order calculus;
Theorem A.6 applies the image theorem and identifies its sentence part.
The first-order dependency chain is A.1 $\to$ A.2 $\to$ A.3 $\to$ A.5,
together with A.4 $\to$ A.5. A.6 uses A.4 and the three propositional
scheme expressions, without the binding or decidable-equality assumptions.

No lost dependency or unrelated section was identified. The list identities,
substitution laws, distinction between formula and sentence closure, and
restriction-splitting paragraph each perform a mathematical task.
The construction gives a generated theory for each axiom subobject;
it does not construct an object classifying all theories.

\newpage
\section*{Implemented corrections}
\textbf{R1: explicit Boolean morphisms (A.1).}
The phrase ``finite coproduct tables'' has been replaced by the actual
copairings. For the injections $\mathsf f,\mathsf t:1\to\Two=1+1$,
distributivity gives
\[
\gamma=[\langle\mathsf f,\mathsf f\rangle,
         \langle\mathsf f,\mathsf t\rangle,
         \langle\mathsf t,\mathsf f\rangle,
         \langle\mathsf t,\mathsf t\rangle]:
         1+1+1+1\xrightarrow{\cong}\Two^2.
\]
The definitions are
\[
\neg_2=[\mathsf t,\mathsf f],\qquad
\wedge_2=[\mathsf f,\mathsf f,\mathsf f,\mathsf t]\gamma^{-1},
\qquad
\vee_2=[\mathsf f,\mathsf t,\mathsf t,\mathsf t]\gamma^{-1}.
\]
Thus $\neg_2:\Two\to\Two$ and
$\wedge_2,\vee_2:\Two^2\to\Two$ have the declared types.
The restrictions to the summands specify the operations; the coproduct
universal property then determines their copairings. The symbol $\gamma$
avoids conflict with the signature's arity bound $d$.

\textbf{R2: the composite in the substitution proof (A.3).}
If $h:I_E\to(X_T,X_F)$ is the initial algebra morphism and
$\pi:(X_T,X_F)\to I_E$ projects to the original expressions, the
endomorphism is $\pi h$. The proof now reads:
\begin{quote}
Composing this algebra morphism with the projections to the original
components gives an endomorphism of the initial $K_E$-algebra, hence the identity.
\end{quote}
This repairs the local typing statement without changing the construction.

\textbf{Explicit ambient category (A.6).}
The fourth agent's local clarification declares the atom object as
$V\in\operatorname{Ob}(\D)$. This matches the first-order declaration
and the author's requirement that objects have an explicit ambient category.

\section*{Writing and terminology findings}
\begin{tabularx}{\textwidth}{@{}L{0.13\textwidth}Y Y@{}}
\toprule
Location & Finding and disposition & Source principle or rationale\\\midrule
A.1 & R1 is resolved by typed copairings. & KLR \S1, points 11 and 26:
explain introduced notation and avoid unnecessary jargon.\\
A.2--A.3 & ``Fold'' and ``carrier'' are established terms. Optional replacements
were considered and declined; no ambiguity was demonstrated. & Audience-sensitive
editorial judgment; technical vocabulary is not prohibited mechanically.\\
A.3 & R2 is resolved; retain the explanation of the equal-name binder branch.
& KLR \S1, points 10--11: connect formulas and provide complementary explanation.\\
A.1--A.6 & Keep the present definition order and proof dependencies.
& KLR \S\S4--5 on contextual definition timing; Ullman, \S30, p.~67, on
types and nearby definitions; \S37, pp.~92--93, on proof order.\\\bottomrule
\end{tabularx}

``Raw named-variable syntax object,'' ``stable under parameters,''
``admissibility,'' and ``sentence theory'' are defined and subsequently used.
Free-variable values are expressly lists of occurrences, retaining repetitions.
No further undefined invented terminology was identified. No new citation,
glossary, reorganization, or numerical realization proof was added.

\newpage
\section*{Independent reports and reconciliation}
\textbf{Agent 1: hostile mathematical review.}
The report checks the types and universal properties in all six sections,
the substitution branches and laws, descent of the inference algebra to its
regular image, leastness, sentence closure, and the partial-inverse condition.
It also checks the relevant realization interfaces in manuscript \S\S4, 8,
and 9. Its findings in the required audit categories are:
\begin{itemize}
\item[(A)] No confirmed hard blocker.
\item[(B)] No established major weakness in the audited construction.
\item[(C)] Both prior local defects resolved; no new definite minor defect.
\item[(D)] No confirmed inconsistency between the appendix copies or the checked
application interfaces.
\item[(E)] No unmet mathematical correction condition supported by this pass.
\item[(F)] Original-source adequacy of the Hilbert presentation, unrelated
external results, novelty, and formal verification were not repeated or established.
\end{itemize}
Its final publication recommendation is unknown outside this bounded audit;
the checked findings supply no ground for rejecting the revised construction.

\textbf{Agent 2: mathematical-writing review.}
The report supplies an actual outline, a source-grounded KLR comparison,
a definition and notation audit, exact optional replacements, and a coverage
record. It finds no further necessary writing change. The source advice is
distinguished from the reviewer's categorical judgments and vocabulary preferences.

\textbf{Agent 3: six-section outline.}
The report assigns each section its function, constructions, principal results,
and prerequisites. It finds no missing dependency or unrelated section and
recommends retaining the existing order.

\textbf{Agent 4: reconciliation and editing.}
The fourth agent read all three new reports, checked their findings against
the appendix, and retained R1 and R2. It declined the optional changes to
``fold'' and ``carrier,'' made the explicit ambient-category declaration
in A.6, and synchronized the standalone and integrated sources.

\section*{Coverage, verification, and limits}
The hostile and outline reports identify the reviewed 598-line source snapshot
by SHA-256. The final edit is recorded separately in the fourth report.
The reviewers reread the complete appendix in this pass; the hostile reviewer
also inspected all seven pages of the corrected PDF. Earlier complete readings
of the manuscript and KLR report are distinguished from this pass's targeted
rereading. The writing review checked the cited KLR advice in the supplied
Stanford report. The four mathematical bibliography sources were not freshly
verified in this pass.

The final standalone appendix and assembled manuscript were compiled, checked
for unresolved references and layout warnings, and inspected at the revised
pages. Their mathematical bodies agree apart from the integrated appendix's
additional compatibility label. The accepted abstract and the other mathematical
sections were preserved. The manuscript version is 14. No Lean verification
is claimed.

\textbf{Review source.} D.~E.~Knuth, T.~Larrabee, and P.~M.~Roberts,
\emph{Mathematical Writing}, Stanford report.
\href{https://www-cs-faculty.stanford.edu/~knuth/papers/cs1193.pdf}{Official report}.
Page citations above use printed pages. The full individual reports and
reconciliation are in
\href{https://github.com/flengyel/CategoricalEssentialIncompleteness/tree/main/Audit/2026-10-31-appendix-v14}{the version 14 audit directory}.
\end{document}
